\chapter{Building Block View (Static View)}

\section{Overview}

\section{Configuration}

LNM configuration is done the same way as for standard BSW modules. A parameter definition file is included with the module which allows 
the users to configure the available parameters.

\subsection{Included Containers}

\input{generated/tables/containers.tex}

\subsection{LNM_General}
\label{sec:LNM_General}

\ConfigurationParameter{LNM_MainFunction_Period}
{Period in seconds of LNM_MainFunction runnable}
{1..1}
{ECUC-FLOAT-PARAM-DEF}
{0.001..0.2}
{0.1}


\ConfigurationParameter{LNM_ScheduleManagement_Period}
{Period in seconds of Schedule Management runnable}
{1..1}
{ECUC-FLOAT-PARAM-DEF}
{0.001..0.2}
{0.1}


\ConfigurationParameter{LNM_NetworkStatus_Period}
{Period in seconds of Network Status runnable}
{1..1}
{ECUC-FLOAT-PARAM-DEF}
{0.001..0.2}
{0.1}


\ConfigurationParameter{LNM_NetworkFaultManager_Period}
{Period in seconds of Network Fault Manager runnable}
{1..1}
{ECUC-FLOAT-PARAM-DEF}
{0.001..0.2}
{0.1}


\ConfigurationParameter{LNM_PermitExecution_Mask}
{Bit mask that must be checked with PermitExecution to allow execution of SUM_LNM functions}
{1..1}
{ECUC-INTEGER-PARAM-DEF}
{0..255}
{254}


\ConfigurationParameter{LNM_WupTxRetries_Commander}
{Init value configuration for the calibration parameter LNM_WupTxRetries_Commander}
{1..1}
{ECUC-INTEGER-PARAM-DEF}
{0..255}
{10}


\ConfigurationParameter{CyberFill_Supported}
{Enable or Disable CyberFill feature}
{1..1}
{ECUC-BOOLEAN-PARAM-DEF}
{NA}
{false}


\ConfigurationParameter{LNM_CyberFillValue}
{Defines the CyberFill value that must be provided by SUM_LNM. This parameter is mandatory if CyberFill functionality is enabled}
{1..1}
{ECUC-INTEGER-PARAM-DEF}
{0..255}
{0}


\ConfigurationParameter{LNM_TWuDelay_Commander}
{Defines the time delay in ms before a table is scheduled during a commander wakeup }
{1..1}
{ECUC-INTEGER-PARAM-DEF}
{0..255}
{100}


\ConfigurationParameter{LNM_TWuDelay_Responder}
{Defines the time delay in ms before a table is scheduled during a responder wakeup }
{1..1}
{ECUC-INTEGER-PARAM-DEF}
{0..255}
{100}


\ConfigurationParameter{LNM_DID_F1BE_Reference}
{Reference to the Dcm configuration of DID F1BE}
{1..1}
{ECUC-REFERENCE-DEF}
{N/A}
{N/A}


\ConfigurationParameter{LnmGenerateInitRunnable}
{Enable or Disable the Runnable Init feature.}
{1..1}
{ECUC-BOOLEAN-PARAM-DEF}
{NA}
{false}

\ConfigurationParameter{LnmInitFuncPeriodicity}
{Specifies time in ms for the periodicity of the runnable Runnable_LNM_Init.}
{0..1}
{ECUC-FLOAT-PARAM-DEF}
{1.0..5000.0}
{1000.0}

\Requirement{510}
{

\ConfigurationParameter{Sum_Lnm_ConfGen_Enabled}
{If false enables the generation of Sum_Lnm_SWCD.arxml, if true the ConfGen will generate that file.}
{1..1}
{ECUC-BOOLEAN-PARAM-DEF}
{NA}
{false}

}
{\Satisfies{SRS_LNM 468}}


\ConfigurationParameter{Sum_Lnm_ScheduleChannelUserFuncs_Enabled}
{Enable the use of user functions for changing the Schedule Table and Channel State for the LIN network.}
{1..1}
{ECUC-BOOLEAN-PARAM-DEF}
{NA}
{false}


\ConfigurationParameterExt{LNM_C_DTC_LostCommWithLinResp_Enabled}
{Initial value for the LNM calibration parameter controlling whether to enable the Lost Communication with LIN Responder Device DTC. Setting this parameter to true will let LNM aggregate the DTC.}
{1..1}
{ECUC-BOOLEAN-PARAM-DEF}
{NA}
{false}

\subsection{LNM_Network}


\ConfigurationParameter{LNM_ActivationSignalType}
{Defines the source of activation signal}
{0..1}
{ECUC-ENUMERATION-PARAM-DEF}
{Com, SWC, Rule}
{SWC}


\ConfigurationParameter{LNM_ActivationMask}
{Defines the bit mask of relevant bits for activation. Bits set to 1 are not relevant for activation}
{0..1}
{ECUC-INTEGER-PARAM-DEF}
{0..255}
{254}


\ConfigurationParameter{LNM_DeactivationMask}
{Defines the bit mask of relevant bits for deactivation. Bits set to 1 are not relevant for deactivation}
{0..1}
{ECUC-INTEGER-PARAM-DEF}
{0..255}
{254}


\ConfigurationParameter{LNM_ChannelEnable_InitValue}
{Define the initial value for the LNM calibration parameter, to enable or disable the MSCL for the selected channel}
{0..1}
{ECUC-BOOLEAN-PARAM-DEF}
{NA}
{true}


\ConfigurationParameter{LNM_MSCL_Strategy}
{Defines the strategy used for MSCL}
{1..1}
{ECUC-ENUMERATION-PARAM-DEF}
{Calibration,Evaluation}
{Calibration}


\ConfigurationParameter{LNM_Network_ID}
{Defines the ID of the LIN network. It is used to select the correct byte position in DID F1BE}
{1..1}
{ECUC-INTEGER-PARAM-DEF}
{1..16}
{1}


\ConfigurationParameter{LNM_C_DTC_LINBusOff}
{Defines whether DTC for LIN Bus off for the given LIN network shall be enabled or not.}
{1..1}
{ECUC-BOOLEAN-PARAM-DEF}
{NA}
{false}


\ConfigurationParameter{Sum_LNM_LinSMNetworkId_ref}
{Reference to the unique LinSM network index}
{0..1}
{ECUC-REFERENCE-DEF}
{N/A}
{N/A}


\ConfigurationParameter{Sum_LNM_Rule_ref}
{Reference to the the activation rule}
{0..1}
{ECUC-REFERENCE-DEF}
{N/A}
{N/A}

\ConfigurationParameter{LNM_C_PreSleepDelay}
{Calibration parameter for Time in ms, indicates the delay from the instance the schedule switching criteria is satisfied to transition from a RUN_ONCE or a RUN_CONTINUOUS FULL_COMM state to a NO_COMM state, until the transition to the NULL schedule occurs when the SLEEP command is issued.}
{1..1}
{ECUC-INTEGER-PARAM-DEF}
{0..4000}
{25}

\subsubsection{LNM_Responders}
\label{subsubsec:LNM_Responders}

\ConfigurationParameter{LNM_C_ResponderPresent_InitValue}
{Define the initial value for the LNM calibration parameter if MSCL strategy is via Calibration}
{0..1}
{ECUC-BOOLEAN-PARAM-DEF}
{NA}
{true}


\ConfigurationParameter{LNM_LossOfComm_withDevice_Timer}
{Defines the amount of time [ms] before set Loss of Communication with Device}
{1..1}
{ECUC-INTEGER-PARAM-DEF}
{0..65535}
{0}


\ConfigurationParameter{LNM_Internal_Electronic_Failure_Supported}
{Define if the responder supports Internal Electronic Failure}
{1..1}
{ECUC-BOOLEAN-PARAM-DEF}
{NA}
{false}


\ConfigurationParameter{LNM_IEF_APINFO2_Supported}
{If Internal Electronic Failure DTC is supported, select if APINFO2 (Reset Status) to detect fault condition is supported}
{0..1}
{ECUC-BOOLEAN-PARAM-DEF}
{NA}
{false}		


\ConfigurationParameter{LNM_IEF_Reset_Threshold}
{If APINFO2 is supported to set Internal Electronic Failure DTC, set the threshold of reset notification that are accepted before setting the DTC}
{0..1}
{ECUC-INTEGER-PARAM-DEF}
{0..255}
{3}


\ConfigurationParameter{LNM_IEF_APINFO4_Supported}
{If Internal Electronic Failure DTC is supported, select if APINFO4 (Application Status) to detect fault condition is supported}
{0..1}
{ECUC-BOOLEAN-PARAM-DEF}
{NA}
{false}


\ConfigurationParameter{LNM_Responder_Device_Not_Configured_Supported}
{Select if the current responder supports the LIN Responder Device Not Configured DTC. It requires APINFO3 signal support }
{1..1}
{ECUC-BOOLEAN-PARAM-DEF}
{NA}
{false}


\ConfigurationParameter{LNM_Config_Retries}
{Number of retries for configuration of the responder before Responder Device Not Configured is set}
{0..1}
{ECUC-INTEGER-PARAM-DEF}
{0..7}
{3}		


\ConfigurationParameter{LNM_MSCL_BitIndex}
{Defines the bit index in MSCL for the selected Responder. MSCL Bit 7 of BYTE 2 is reserved for the Config Bit}
{1..1}
{ECUC-INTEGER-PARAM-DEF}
{0..7}
{0}


\ConfigurationParameter{LNM_MSCL_ByteIndex}
{Defines the MSCL Byte of the selected Responder}
{1..1}
{ECUC-ENUMERATION-PARAM-DEF}
{MSCL_BYTE_1,MSCL_BYTE_2}
{MSCL_BYTE_1}


\subsubsection{LNM_ScheduleStates}

\ConfigurationParameter{LNM_ScheduleRef}
{Reference to the schedule table associated to the corresponding state}
{0..1}
{ECUC-REFERENCE-DEF}
{N/A}
{N/A}


\subsubsection{LNM_ScheduleTransitions}


\ConfigurationParameter{LNM_StartStateRef}
{Reference to the start state of the corresponding transition}
{0..1}
{ECUC-REFERENCE-DEF}
{N/A}
{N/A}


\ConfigurationParameter{LNM_EndStateRef}
{Reference to the end state of the corresponding transition}
{0..1}
{ECUC-REFERENCE-DEF}
{N/A}
{N/A}


\ConfigurationParameter{LNM_TransitionRuleRef}
{Reference to the rule that is implementing the logic of the selected transition}
{0..1}
{ECUC-REFERENCE-DEF}
{N/A}
{N/A}


\subsection{LNM_Arbitration}
\label{subsec:LNM_Arbitration}

\subsubsection{LNM_ScheduleSignal}


\ConfigurationParameter{LNM_ScheduleSignalSource}
{Defines the source of schedule signal}
{0..1}
{ECUC-ENUMERATION-PARAM-DEF}
{Com, SWC, VPM}
{SWC}


\ConfigurationParameter{LNM_ScheduleSignalType}
{Defines the type of schedule signal}
{0..1}
{ECUC-ENUMERATION-PARAM-DEF}
{boolean, uint8, uint16}
{uint8}


\subsubsection{LNM_Condition}


\ConfigurationParameter{LNM_CompareOperator}
{Defines the logic operation that is performed on input signal}
{0..1}
{ECUC-ENUMERATION-PARAM-DEF}
{==, !=, <, >, <=, >=}
{==}


\ConfigurationParameter{LNM_CompareValue}
{Defines the value to which the input signal is compared. In case of boolean variables it is assumed that: TRUE == 1 and BOOLEAN == 0}
{0..1}
{ECUC-INTEGER-PARAM-DEF}
{0..65535}
{0}


\ConfigurationParameter{LNM_CompareVPMValue}
{Defines the value to which the vehicle power mode is compared.}
{0..1}
{ECUC-ENUMERATION-PARAM-DEF}
{OFF, ACCESSORY, RUN, STARTREQUEST, PROPULSION, TRANSPORT}
{N/A}


\ConfigurationParameter{LNM_InputRef}
{Reference to the input signal of the selected condition}
{0..1}
{ECUC-REFERENCE-DEF}
{N/A}
{N/A}


\subsubsection{LNM_Rule}


\ConfigurationParameter{LNM_LogicOperator}
{Defines the logic operation that relates all the respective referenced conditions and rules. Logic AND (\&), Logic OR (|)}
{0..1}
{ECUC-ENUMERATION-PARAM-DEF}
{\&, |}
{\&}


\ConfigurationParameter{LNM_ConditionRef}
{Reference to basic LNM_Condition or LNM_Rule to build complex logic operations. IMPORTANT: Since LNM_Rules can reference others LNM_Rules, recursion must be carefully avoided at configuration time}
{0..Inf}
{ECUC-REFERENCE-DEF}
{N/A}
{N/A}


